\section{引言}\label{sec:introduction}
图表示学习将节点属性与关系结构映射为可用于分类等任务的向量。与独立同分布样本不同，图中节点的预测依赖其自身特征和邻域信息；但获取节点类别标签通常需要额外成本，少量标签难以直接覆盖图的复杂结构。图神经网络可通过邻域聚合联合利用结构与属性，图自监督学习则进一步从未标注数据中构造训练信号，再将所得表示用于下游任务。

GraphMAE是Hou等发表于KDD 2022的生成式图自监督方法\cite{hou2022graphmae}。它关注一个具体问题：如何使图自编码器学到有用表示，而不是仅完成容易的输入复制或结构拟合？其设计以节点特征为重建目标，通过输入掩码、图神经网络解码、解码前重掩码及缩放余弦误差形成自监督任务。结构仍用于消息传递，但不被作为该重建损失的预测对象。本文选择该论文开展课程级部分复现，范围限定为Cora上的节点分类，不覆盖原论文的全部数据集和任务。

本文的工作包括两部分。首先，固定官方代码版本、数据处理和超参数，完成五种子复现，并将验证集选模与测试集最终评价分离。其次，提出并检验度数感知掩码假设：高度数节点具有较多邻域上下文，提高其掩码频率是否有助于表示学习？为此只调整掩码节点选择分布，在固定预算下对三个正向系数和均匀控制组进行配对实验。GraphMAE本身的架构与损失贡献属于原作者；本文的自主工作是采样适配、随机数流控制及相应实验分析。实验并未获得平均分类性能提升，因此本文保留负结果，并分析其证据边界。
